\noindent\textbf{Elección del procedimiento.} La solución exterior $u_\theta=\Omega R^2/r$ da $\tau_{r\theta}(R)=-2\mu\Omega$. Su par por longitud tiene módulo $T\prime=4\pi\mu\Omega R^2$ y $P\prime=T\prime\Omega$. Del balance radial $dp/dr=\rho u_\theta^2/r$ se obtiene $p_\infty-p(R)=\rho\Omega^2R^2/2$.

\begin{enumerate}

\item Comparar la escala de rotación con viscosidad.

\[\fitmath{\mathrm{Re}=\frac{\rho\cdot W\cdot {R}^{2}}{\mu}}\]

\[\fitmath{\mathrm{Re}=\frac{(970)\cdot (50)\cdot {(0.02)}^{2}}{(0.5)}\approx 38.8\,}\]

\item La presión en la pared es menor que lejos del cilindro.

\[\fitmath{\mathrm{dp}=\frac{\rho\cdot {W}^{2}\cdot {R}^{2}}{2}}\]

\[\fitmath{\mathrm{dp}=\frac{(970)\cdot {(50)}^{2}\cdot {(0.02)}^{2}}{2}\approx 485\,\mathrm{Pa}}\]

\item Integrar esfuerzo por brazo por unidad de longitud axial.

\[\fitmath{T\prime=4\cdot \pi\cdot \mu\cdot W\cdot {R}^{2}}\]

\[\fitmath{T\prime=4\cdot \pi\cdot (0.5)\cdot (50)\cdot {(0.02)}^{2}\approx 0.1256637\,\mathrm{N}}\]

\item Calcular la potencia necesaria por unidad de longitud.

\[\fitmath{P\prime=T\prime\cdot W}\]

\[\fitmath{P\prime=(0.1256637061)\cdot (50)\approx 6.283185\,\mathrm{W/m}}\]

\end{enumerate}
